\chapter{Subcontractor, Supply Chain and Procurement Risk Assessment}
\label{chap:Subcontractor, Supply Chain and Procurement Risk Assessment}

%---------------------------- SECTION ----------------------------
\section{Why this assessment matters in construction}

\paragraph{The modern construction project is not delivered by a single organisation. It is delivered by a complex, layered supply chain in which the principal contractor retains overall CDM responsibility while the actual physical work is performed by a constellation of specialist subcontractors, sub-tier suppliers, material and plant hirers, logistics providers, and visiting specialists. The health and safety performance of the overall project is therefore directly dependent on the health and safety performance of every link in that supply chain --- and the principal contractor who selects, appoints, and manages those supply chain relationships is legally accountable for ensuring that each link is strong enough to carry the load.}

\paragraph{The pre-qualification questionnaire (PQQ) and the SSIP (Safety Schemes in Procurement) framework represent the industry's primary mechanism for filtering out supply chain organisations that lack basic health and safety competence before they are appointed. SSIP-accredited schemes --- including CHAS (Contractors Health and Safety Assessment Scheme), SafeContractor, and Constructionline Gold --- assess subcontractors against a core criteria set that covers health and safety policy, risk assessment capability, accident history, insurance, and CDM awareness. SSIP accreditation provides a starting point for competence assurance, not an end point; the principal contractor's procurement process must go beyond SSIP to assess project-specific capability for the tasks to be undertaken.}

\paragraph{Risk Assessment and Method Statements (RAMS) are the principal contractor's primary tool for ensuring that subcontractors have planned their work safely before committing operatives to the hazardous activity. The RAMS review and approval process --- through which the principal contractor's safety team reviews the subcontractor's RAMS before any work begins, identifies gaps or inadequacies, requires revisions, and formally approves the document --- is a critical control in the construction health and safety management system. A RAMS that is rubber-stamped by a busy commercial manager who has not read it, or that is accepted without review because the subcontractor has SSIP accreditation, provides no risk management benefit and may create a false sense of assurance that the activity has been planned safely.}

\paragraph{Modern slavery represents a supply chain risk that construction companies of a significant scale are required by law to assess and manage. The Modern Slavery Act 2015 requires organisations with an annual turnover of \pounds36\,million or more to publish an annual slavery and human trafficking statement describing the steps they are taking to ensure that modern slavery is not occurring in their operations or supply chains. The construction industry --- with its complex, multi-tier supply chains, its significant migrant workforce, its pattern of cash-in-hand labour supply, and its exposure to the exploitation networks that operate in the labour supply sector --- is identified by the National Crime Agency as a high-risk sector for modern slavery and labour exploitation.}

\barquote{``The principal contractor who delegates the physical work to subcontractors does not delegate the CDM duty to ensure that the work is carried out safely. The supply chain is the vehicle through which health and safety is delivered; and if the vehicle is not roadworthy, the duty holder who put it on the road shares responsibility for the crash.''}

\example{Supply Chain Scenario}{A principal contractor appoints a groundworks subcontractor for the foundations and drainage package of a residential development. The subcontractor holds current SSIP accreditation and has submitted a RAMS for the excavation work. The principal contractor's safety manager reviews the RAMS and approves it without modification. Three weeks into groundworks, an HSE inspection reveals that the RAMS makes no reference to the CAT survey requirement for the excavation area, the trench box specification is insufficient for the soil conditions identified in the pre-construction information, and the RAMS was produced using a generic template with no site-specific risk assessment. Work is stopped. The HSE's investigation notes that the principal contractor's RAMS approval process did not include any check of RAMS adequacy against the site-specific pre-construction information.}

%---------------------------- SECTION ----------------------------
\section{Legal and professional framework}

\paragraph{The principal contractor's responsibility for the health and safety of subcontractors on site is embedded in CDM 2015 Regulations 12 to 14, which require the principal contractor to plan, manage, monitor, and coordinate health and safety in the construction phase --- including the activities of all contractors on site. The duty to ensure that contractors are competent is found in CDM 2015 Regulations 8 and 9, which require the appointing duty holder to satisfy themselves that those appointed have the skills, knowledge, experience, and organisational capability to perform their CDM role. The general duty under Section 3 of the Health and Safety at Work etc.\ Act 1974 requires the principal contractor to ensure, so far as is reasonably practicable, that persons not in their employment who may be affected by the conduct of the construction work are not exposed to risks to their health and safety.}

\paragraph{The Modern Slavery Act 2015, the Late Payment of Commercial Debts (Interest) Act 1998, and the CLOCS (Construction Logistics and Cyclist Safety) and FORS (Fleet Operator Recognition Scheme) standards represent a broader supply chain governance framework that goes beyond health and safety into ethical, financial, and logistical compliance. Together they reflect the reality that the supply chain relationship in construction is not merely a commercial transaction but a shared responsibility for the welfare of the workers, the communities, and the environment affected by the construction activity.}

\begin{itemize}
  \bitem{CDM 2015 Regulations 8, 9, and 12--14}{Regulation 8 prohibits the appointment of a contractor or designer unless the appointing duty holder is satisfied that the appointee has the skills, knowledge, experience, and organisational capability to perform their role. Regulations 12--14 place the construction phase management, coordination, and monitoring duties on the principal contractor. The principal contractor must ensure that contractors comply with the Construction Phase Plan; that there is cooperation between contractors; and that the site rules are communicated to and observed by all.}
  \bitem{SSIP (Safety Schemes in Procurement) framework}{SSIP is an umbrella body whose member schemes (CHAS, SafeContractor, Constructionline Gold, Acclaim, PAS 91, and others) all assess subcontractors against a common core criteria set. SSIP accreditation provides evidence of basic health and safety competence and is widely accepted by principal contractors as part of their supplier approval process. SSIP member schemes allow mutual recognition of each other's assessments, reducing the burden of multiple pre-qualification submissions on subcontractors.}
  \bitem{RAMS review and approval process}{The principal contractor must establish a documented process for receiving, reviewing, approving, and communicating RAMS from subcontractors before any work commences. The review must be conducted by a person competent to assess the adequacy of the RAMS for the specific activity and site conditions; a RAMS that does not reflect the pre-construction information, the site conditions, or the specific hazards of the activity must be returned for revision before approval.}
  \bitem{Modern Slavery Act 2015}{Requires organisations with annual turnover of \pounds36\,million or more to publish an annual slavery and human trafficking statement. In practice, all construction organisations should assess their supply chains for indicators of modern slavery including: workers who are transported in and out of site by a labour provider; workers who share accommodation provided by their employer; workers who cannot speak to the principal contractor's staff independently; and workers who appear to be having wages deducted for accommodation, transport, or tools.}
  \bitem{Late Payment of Commercial Debts (Interest) Act 1998 and the Construction Act 1996}{Supply chain payment risk is a significant contributing factor to health and safety risk in construction: subcontractors under severe financial pressure cut corners on safety; workers who have not been paid may not attend site, disrupting the programme and creating gaps in the supervision structure. The principal contractor's procurement process should assess the financial stability of subcontractors as part of the pre-qualification process.}
  \bitem{CLOCS Standard for Construction Logistics}{The CLOCS (Construction Logistics and Cyclist Safety) standard requires principal contractors and their supply chain logistics providers to manage the safety of vulnerable road users (cyclists, pedestrians, motorcyclists) who interact with delivery vehicles travelling to and from the construction site. CLOCS compliance requires that delivery vehicles meet minimum construction logistics safety standards; that delivery windows and routes are planned to minimise conflict with vulnerable road users; and that all logistics providers are CLOCS-compliant.}
  \bitem{FORS (Fleet Operator Recognition Scheme)}{FORS is a voluntary accreditation scheme for vehicle operators that sets standards for safety, efficiency, and environment. FORS Bronze, Silver, and Gold accreditation levels are increasingly required by principal contractors for logistics providers whose vehicles access their sites. FORS-accredited logistics providers are assessed against standards that address driver licensing, vehicle maintenance, CCTV and sensors on vehicles, and management systems for road safety incidents.}
  \bitem{Checking insurance and certificates}{Principal contractors must verify that all subcontractors hold valid employer's liability insurance (minimum \pounds5\,million, typically \pounds10\,million), public liability insurance (minimum \pounds5\,million, typically \pounds10\,million), and any specialist professional indemnity or product liability insurance required for their scope. Insurance certificates must be checked at pre-qualification and again at appointment; expired insurance is a common supply chain risk gap.}
\end{itemize}

%---------------------------- SECTION ----------------------------
\section{Conducting the assessment using the five-step method}

\subsection{Step 1 --- Identify Hazards}
\paragraph{Supply chain and procurement hazards arise from failures in the selection, assessment, and management of subcontractors. Key hazards include: appointment of a subcontractor without adequate health and safety competence for the specific tasks to be undertaken; appointment of a subcontractor with an inadequate or falsified SSIP certificate; approval of RAMS that do not adequately address the hazards of the activity or the site conditions; failure to monitor subcontractor compliance with their RAMS during work execution; appointment of a logistics provider whose vehicles do not meet the CLOCS standard, creating road safety risk for cyclists and pedestrians; failure to identify indicators of modern slavery or labour exploitation in the supply chain; failure to verify insurance currency at appointment, leaving the principal contractor without recourse in the event of a supply-chain-caused incident; and engagement of visiting specialists (commissioning engineers, manufacturers' representatives, specialist installers) without applying the same RAMS and induction requirements that apply to regular subcontractors.}

\subsection{Step 2 --- Decide Who or What Is Affected}
\paragraph{Workers employed by subcontractors are the primary persons at risk from inadequate supply chain management; they are working on the principal contractor's site, subject to the principal contractor's CDM responsibilities, but their employer's health and safety culture and risk assessment capability may be significantly inferior to the principal contractor's. Members of the public who interact with delivery vehicles are at risk from logistics providers who do not meet CLOCS standards. Future occupants and maintainers of the building are at risk from sub-standard workmanship or non-compliant materials installed by an inadequately assessed supply chain. The principal contractor is at financial and reputational risk from supply chain failures that cause programme delay, insurance gaps, or regulatory prosecution.}

\subsection{Step 3 --- Evaluate Likelihood and Consequence}
\paragraph{The consequence of appointing an incompetent or dishonest subcontractor can be fatal for a site operative. The likelihood of supply chain health and safety failures increases in direct proportion to the pressure on procurement budgets: a subcontractor who wins a package at a price that leaves no margin for adequate safety management will cut corners, and the cut that proves fatal may not be attributable directly to the low tender price but the causal chain will lead back to the procurement decision. The consequence of undetected modern slavery in the supply chain includes criminal liability for the principal contractor under the Modern Slavery Act, reputational damage, and complicity in the exploitation of vulnerable workers.}

\subsection{Step 4 --- Select and Implement Controls}
\paragraph{Controls for supply chain risk must operate across three stages: pre-qualification, appointment, and ongoing monitoring. At pre-qualification: require SSIP accreditation as a minimum threshold; verify SSIP certificate validity directly with the issuing scheme; conduct project-specific competence assessment for the scope of work; verify insurance currency; check accident history and RIDDOR record; verify financial stability; and check for CLOCS/FORS compliance for logistics providers. At appointment: issue a formal purchase order or subcontract that incorporates the site rules, CDM obligations, and RAMS submission requirement; require RAMS submission and approval before work starts. During execution: conduct subcontractor safety inductions; verify CSCS card compliance for all operatives; monitor RAMS compliance through daily site supervision and formal safety inspections; conduct sub-tier supply chain checks for any sub-tier suppliers engaged by the subcontractor without the principal contractor's knowledge; and report and investigate any subcontractor health and safety incident under the principal contractor's incident reporting system.}

\subsection{Step 5 --- Record and Review}
\paragraph{Supply chain records must be maintained throughout the project: pre-qualification records, appointment documentation, RAMS approval records, induction records, CSCS card logs, and incident records. At project close-out, the supply chain performance data --- accidents, near misses, RAMS rejection rates, SSIP compliance issues --- should be fed back into the approved supplier database to inform future procurement decisions. The Modern Slavery Act statement must be reviewed annually and updated to reflect supply chain assessments conducted during the year. SSIP certificates must be re-verified annually for subcontractors with whom the principal contractor has a continuing relationship.}

\example{Five-Step Application}{A principal contractor is assembling the supply chain for a 60-week mixed-use development. Step 1 identifies: two proposed subcontractors on the approved list whose SSIP certificates expire within three months of the planned appointment date; a drainage subcontractor whose PQQ indicates no previous experience of CDM-notifiable projects; the proposed concrete frame subcontractor engages a sub-tier reinforcement supply chain that the principal contractor has not assessed. Step 2 confirms: reinforcement workers are on site under the principal contractor's CDM responsibility but without PQQ or RAMS review. Step 3 rates risk from unapproved sub-tier supply chain as High. Step 4 implements: SSIP re-verification required before appointment; drainage subcontractor required to demonstrate CDM CPP experience with references; sub-tier reinforcement supplier required to complete PQQ and submit RAMS before any work starts; frame subcontractor contractually required to obtain principal contractor approval before engaging any sub-tier supplier. Step 5 specifies: quarterly supply chain audit; annual SSIP re-verification cycle.}

%---------------------------- SECTION ----------------------------
\section{Applying the hierarchy of controls}

\paragraph{The hierarchy of controls for supply chain risk management places the emphasis on pre-qualification and selection controls --- the highest-order interventions that prevent incompetent or dishonest organisations from entering the site in the first place --- ahead of monitoring and enforcement controls that attempt to manage non-compliance after the appointment is made.}

\begin{itemize}
  \bitem{Elimination}{Eliminate supply chain health and safety risk at its source by engaging only SSIP-accredited subcontractors with verified, current certificates; by requiring project-specific competence demonstration for complex or high-risk scopes; and by not awarding packages to the lowest tender price where that price is demonstrably insufficient to fund adequate safety management.}
  \bitem{Substitution}{Where a proposed subcontractor cannot demonstrate adequate competence for the specific scope, substitute an alternative supplier from the approved list; where CLOCS-compliant logistics providers are available for the delivery schedule required, substitute non-compliant providers.}
  \bitem{Engineering controls}{Implement an electronic supplier management system that flags expiring SSIP certificates, insurance renewals, and RAMS approval status; use a RAMS management platform that requires electronic approval before the start of any new activity.}
  \bitem{Administrative controls}{Implement a documented PQQ process; establish SSIP accreditation as a minimum pre-qualification condition; implement a RAMS submission and approval procedure; conduct subcontractor inductions; verify CSCS card compliance at each new operative's first site attendance; conduct monthly supply chain safety performance reviews; implement a sub-tier supply chain approval requirement in all subcontracts; include modern slavery awareness in site inductions.}
  \bitem{PPE}{PPE in the supply chain context represents the contractual protections that are the last line of defence: the subcontract terms that require SSIP compliance, RAMS submission, and CDM adherence as conditions of payment; the insurance requirements that ensure financial recourse in the event of a supply-chain-caused incident; and the indemnity clauses that protect the principal contractor from the consequences of the subcontractor's negligence.}
\end{itemize}

%---------------------------- SECTION ----------------------------
\section{Cadence of assessment and review triggers}

\paragraph{Supply chain risk management is a continuous process throughout the project lifecycle, with fixed review points at each stage of the procurement and monitoring cycle and reactive triggers for non-compliance events.}

\begin{itemize}
  \bitem{Pre-appointment PQQ review}{Before any subcontractor is appointed, the PQQ must be reviewed and approved. SSIP certificates must be verified directly with the scheme. Insurance certificates must be checked for currency and adequacy. This is a mandatory gateway; no appointment should proceed without completion.}
  \bitem{Before each new scope of work starts}{RAMS must be submitted, reviewed, and approved before any new scope of work begins, regardless of whether the subcontractor is already on site undertaking a different scope. A new scope of work is a new risk event requiring fresh assessment.}
  \bitem{Monthly supply chain performance review}{Accident statistics, near-miss reports, RAMS rejection rates, and SSIP compliance status for all active subcontractors should be reviewed monthly by the principal contractor's safety and commercial teams.}
  \bitem{On engagement of a new sub-tier supplier}{Whenever a subcontractor proposes to engage a sub-tier supplier, the principal contractor must apply the same PQQ and SSIP verification process to the sub-tier supplier as to a direct appointment.}
  \bitem{After any supply-chain-caused incident}{Any accident, near miss, or regulatory enforcement involving a subcontractor must trigger a review of the subcontractor's RAMS, induction records, and SSIP status, and a decision on whether the subcontractor should remain on the approved supplier list.}
  \bitem{Annual approved supplier list review}{The approved supplier list should be reviewed annually to remove lapsed SSIP certificates, poor performers, and organisations whose insurance or financial position has materially deteriorated.}
\end{itemize}

%---------------------------- SECTION ----------------------------
\section{Common failure modes}

\begin{itemize}
  \bitem{SSIP certificate accepted without verification}{SSIP certificates are available for purchase on the secondary market and are routinely submitted by subcontractors who hold them for other schemes but are presenting the certificate for a scope they are not assessed for. The principal contractor who accepts an SSIP certificate at face value without verifying its currency and scope coverage directly with the issuing scheme has not discharged the competence assessment duty under CDM 2015.}
  \bitem{RAMS approved without substantive review}{The RAMS approval process is one of the most consistently underfunded and under-resourced activities in principal contractor safety management. Safety managers under workload pressure approve RAMS on the basis of their format (does it have the right sections?) rather than their content (does it address the actual hazards of the activity on this site?). A RAMS that says ``all operatives will wear PPE'' without specifying what PPE, why, and when has not been reviewed.}
  \bitem{Sub-tier supply chain not assessed}{Principal contractors who conduct thorough PQQ processes for their direct subcontractors frequently have no visibility of the sub-tier supply chain engaged by those subcontractors. A tier-one subcontractor who sub-contracts specialist elements to a sole trader with no formal health and safety management, no insurance, and no SSIP accreditation has introduced a serious health and safety risk to the site without the principal contractor's knowledge.}
  \bitem{Modern slavery indicators not recognised or reported}{Indicators of modern slavery in the construction supply chain --- workers who cannot speak to site staff independently, workers whose wages appear to be going directly to their labour provider, workers who are clearly exhausted or malnourished, workers who are transported to site in a controlled manner --- are frequently observed and not reported, either because site staff do not recognise them as modern slavery indicators or because they do not know who to report them to.}
  \bitem{Insurance not re-verified at appointment or during the project}{Subcontractors whose insurance was verified at pre-qualification and then never re-verified during the project are a common gap. Annual insurance renewals may result in a change of insurer, a change in policy terms, or a lapse in coverage that the principal contractor does not become aware of until an incident occurs.}
\end{itemize}

\example{Failure Pattern}{A principal contractor appoints a specialist timber frame subcontractor for a residential development. The subcontractor submitted a valid SSIP certificate at pre-qualification. Two months into the project, an operative falls from the timber frame and sustains a specified injury. The HSE investigation reveals: the SSIP certificate covered the subcontractor for general building work but not for timber frame erection above ground floor; the RAMS for the timber frame erection had been submitted and approved by the principal contractor's commercial manager (not the safety manager) without review of the working at height controls; and the timber frame team included two workers engaged by a sub-tier labour agency that the principal contractor had not assessed. Prosecution under CDM 2015 Regulation 8 (failure to satisfy competence) follows.}

%---------------------------- CHAPTER SUMMARY ----------------------------
\newpage
\section{Chapter Summary}

\begin{table}[H]
\centering
\begin{tabularx}{\textwidth}{>{\raggedright\arraybackslash}p{3.2cm}X}
\toprule
\textbf{Assessment Element} & \textbf{Operational Requirement} \\
\midrule
Legal/Professional Basis & CDM 2015 Regulations 8, 9, and 12--14; Health and Safety at Work etc.\ Act 1974 Section 3; Modern Slavery Act 2015; SSIP framework; CLOCS; FORS. \\
Method & PQQ and SSIP verification; project-specific competence assessment; RAMS submission, review, and approval; sub-tier supply chain assessment; insurance verification; CLOCS/FORS compliance check. \\
Controls & SSIP accreditation as minimum pre-qualification condition; direct SSIP certificate verification; RAMS approval by competent reviewer before work starts; sub-tier supplier approval requirement in subcontracts; modern slavery awareness in inductions; insurance re-verification at appointment. \\
Cadence & Pre-appointment PQQ; pre-work RAMS approval; monthly performance review; on sub-tier engagement; after any supply-chain incident; annual approved supplier list review. \\
Assurance & Supply chain safety audit; RAMS rejection rate analysis; accident and near-miss data by subcontractor; annual modern slavery statement review; SSIP certificate currency dashboard. \\
\bottomrule
\end{tabularx}
\end{table}

\begin{itemize}
  \bitem{Key takeaway 1}{SSIP accreditation is a starting point for supply chain competence assurance, not an end point; the principal contractor must conduct project-specific competence assessment for the scope of work to be undertaken, not rely solely on the SSIP certificate.}
  \bitem{Key takeaway 2}{RAMS review and approval must be conducted by a person competent to assess the adequacy of the RAMS for the specific activity and site conditions; rubber-stamping RAMS without substantive review is a common and legally significant control failure.}
  \bitem{Key takeaway 3}{Sub-tier supply chain management is a CDM responsibility of the principal contractor; any organisation working on the site, at whatever tier, must be assessed for competence and covered by the principal contractor's health and safety management system.}
  \bitem{Key takeaway 4}{Modern slavery indicators in the construction supply chain are well-documented and should be known to every person with supply chain management responsibility; unrecognised or unreported indicators represent a legal risk under the Modern Slavery Act and a moral failure.}
  \bitem{Key takeaway 5}{Insurance verification is a procurement gateway that must be repeated at appointment and at annual renewal; lapsed or inadequate insurance discovered after an incident is a financial risk that appropriate procurement processes would have prevented.}
\end{itemize}

\barquote{In Chapter~\ref{chap:Incident Investigation, Reporting and Lessons Learned Risk Assessment}, we examine the legal requirements and best practice for investigating construction incidents, reporting under RIDDOR, and embedding the lessons learned into the risk management cycle.}
